\section{第4課 部屋に　机と　椅子が　ありません}

\subsection{基本課文}

\begin{enumerate}
    \item 部屋に　机と　椅子が　ありません。
    \item 机の　上に猫がいます。
    \item 売店は　駅の　外に　あります。
    \item 吉田さんは　庭に　います。
\end{enumerate}

\begin{itemize}
    \item 甲：その　箱の　中に　何が　ありますか。
    \item 乙：時計と　眼鏡が　あります。
\end{itemize}


\begin{itemize}
   \item 甲：部屋に　誰が　いますか。
   \item 乙：誰もいません。
\end{itemize}

\begin{itemize}
   \item 甲：小野さんの　家は　どこに　ありますか。
   \item 乙：横浜に　あります。
\end{itemize}

\begin{itemize}
   \item 甲：あそこに　犬が　いますね。　
   \item 乙：ええ、　私の　犬です。
\end{itemize}


\subsection{语法}

\begin{itemize}
    \item 〜に　〜が　あります、います。
    \item 〜は　〜に　あります、います。
    \item 上、中、下、前、後ろ、隣、外。
    \item 疑问词　+　も　+　动(否定)：全面否定。
\end{itemize}

\subsection{表达}

\begin{itemize}
    \item 上。
    \item ご家族、ご兄弟、ご両親
    \item 兄弟
    \item 一人暮らし
    \item 建筑设施
\end{itemize}

\subsection{应用课文:会社の場所}

李：小野さん、会社は　どこに　ありますか。
小野：ええと、ここです。
李：近くに　駅が　ありますか。
小野：ええ、JRと　地下鉄の駅が　あります。JRの駅は　ここです。
李：地下鉄の駅は　ここですね。
小野：ええ、そうです。　JRの　駅の　隣に　地下鉄の　駅があります。

\vspace{1em}

李：小野さんの　家は　どちらですか。
小野：私の　家は　横浜です。
李：ご家族も　横浜ですか。
小野：いいえ、私は　一人暮らしです。
李：ご両親は　どちらですか。
小野：両親は　名古屋に　います。
李：ご兄弟は？
小野：大阪に　妹が　います。

\begin{vocablist}[2]
   \vocab{部屋}
   \vocab{庭}
   \vocab{家}
   \vocab{居間}
   \vocab{冷蔵庫}
   \vocab{壁}
   \vocab{スイッチ}
   \vocab{本棚}
   \vocab{ベッド}
   \vocab{ねこ}
   \vocab{犬}
   \vocab{箱}
   \vocab{眼鏡}
   \vocab{ビデオ}
   \vocab{サッカーボール}
   \vocab{ビール}
   \vocab{ウイスキー}
   \vocab{子供}
   \vocab{兄弟}
   \vocab{両親}
   \vocab{妹}
   \vocab{男}
   \vocab{女}
   \vocab{生徒}
   \vocab{上}
   \vocab{外}
   \vocab{中}
   \vocab{下}
   \vocab{前}
   \vocab{後ろ}
   \vocab{近く}
   \vocab{場所}
   \vocab{教室}
   \vocab{会議室}
   \vocab{花屋}
   \vocab{売店}
   \vocab{駅}
   \vocab{地下鉄}
   \vocab{国鉄}
   \vocab{私鉄}
   \vocab{木}
   \vocab{一人暮らし}
   \vocab{あります}
   \vocab{います}
   \vocab{横浜}
   \vocab{名古屋}
   \vocab{大阪}
   \vocab{ジェーアール}
   \vocab{ええと}
\end{vocablist}
